\section{Introduction}

% CTMC models of point substitution are used widely throughout molecular evolution

% One of the most commonly used models is LG08
% Estimated using Expectation Maximization
% Sufficient statistics of endpoint-conditioned CTMCs: expected substitution counts & dwell times \cite{HolmesRubin2002}

% Despite importance of indels in molecular evolution, models of indels are less widely used in phylogenetics
% Canonical indel model is TKF91/TKF92
% Combines a birth-death process with a CTMC
% Despite some attention, earlier work did not produce closed-form expectations for indel/birth-death processes \cite{Holmes2005,CrawfordSuchard2012,DossEtAl2013,CrawfordMininSuchard2014}

% The challenge of placing multiple sequence alignment on a sound likelihood-phylogenetics footing has been called ``Statistical Alignment''. \cite{Hein2000}
% Leading MSA tools such as PRANK \cite{LoytynojaGoldman2008} and MUSCLE \cite{Edgar} use Statistical Alignment-inspired ideas like iterative refinement and consensus over ensembles of trees
% However, going past inspiration to a rigorous reconciliation of MSA algorithms with phylogenetic models has proceeded more slowly; BAliPhy is a notable contribution \cite{RedelingsSuchard2005}

% In part this is because the models used by Statistical Alignment (TKF91/TKF92) still lack features such as context-dependent indel parameters that have been in tools like CLUSTAL since the 1980s
% This modeling gap has grown even larger now that protein modeling space has grown to include methods capable of capturing nonlocal/epistatic interactions such as Potts models \cite{DCA} and autoregressive neural models \cite{arDCA}

\section{Results}

% We first embarked on a thorough analysis and characterization of the TKF family of models,
% in order to develop machinery to fit and apply more elaborate versions of these models

\subsection{The expected indel count and integrated sequence length are expressible in closed form.}

% These results apply more fundamentally to the linear birth-death process with constant immigration:
% the posterior expectations of the sufficient statistics on an endpoint-conditioned trajectory
% are obtainable from the Fisher score identity, by differentiating the transition probability.
% Since the transition probability of the birth-death process is known in closed form \cite{Kendall1948},
% the sufficient statistics follow.
% This result turns out to be remarkably simple given past effort expended in analyzing this model \cite{Holmes2005,CrawfordSuchard2012,DossEtAl2013,CrawfordMininSuchard2014}

% An issue that arises with the TKF model in practice is that the expected sequence length is $1/(1-\insrate/\delrate)$
% Thus, for long sequences, the insertion rate asymptotically approaches the deletion rate
% This can lead to numerical instabilities when calculating expressions that involve $\delrate-\insrate$ in the denominator
% As part of our analysis, we systematically derive expressions for relevant TKF

% A direct product of the closed-form expressions for the expected sufficient statistics is fast Expectation Maximization algorithms with exact M-steps
% We earlier derived the form of this algorithm for general point substitution models \cite{HolmesRubin2002}; in the appendix, we give specific forms of the EM algorithm
% for a variety of point substitution models encountered in the molecular evolution literature,
% including the HKY85 and Goldman-Yang models.

% As with the sufficient statistics, the existence of an EM algorithm is of direct relevance to birth-death models,
% e.g. those used to study gene family evolution \cite{Hahn2005}
% In the setting of indel models, it yields an expression that is linear in the expected transition counts
% Thus, in TKF-like models, we can directly convert the Baum-Welch algorithm into rate parameter space

\subsection{Elaborating evolutionary models with hierarchical mixtures increases their explanatory power.}

% Training data/splits: clan-aware PFam
% Model-fitting: Adam, MixDom Baum-Welch

% SiteMix, FragMix, DomMix
% Annabel's graphs

\subsection{TKF-based models have interpretable parameters.}

% Bubbleplots of domain rates e.g. for d3f1
% Interpretation of domain parameters and top-level domains

% Observation that fitting more ragged/gappy alignments leads to higher top-level rates

% Protein structure figures

\subsection{The new evolutionary models produce more accurate sequence alignments than earlier models, and are comparable to leading MSA methods optimized for accuracy}

% Experimental design
% BAliBase & OxBench SP/TC tables
% Observation: adding domains only helps alignment accuracy up to two or three domains

% Observation: training by summing over alignments doesn't help test-set alignment likelihood (but does improve test-set sum-over-alignments likelihood).
% Best alignments from models trained on curated cherries.

% Observation: P(FSA's MSA|FSA's MSA's tree) > P(BaliBase|BaliBase's tree). 

% EXPERIMENT TO TRY: when training, sum over alignments consistent with given matches

% EXPERIMENT TO TRY: NeuralTKF Forward-Backward -> FSA

\subsection{MixDom matches but does not surpass Felsenstein for substitutions with Fitch parsimony for indels}


\section{Discussion}

% TKF92 Baum-Welch enables training of other HMM-based models with (rate * time) parameterizations



\section{Acknowledgements}

% Funding

% Claude


